\subsection{Superimposed integrals of positive functions}

For the rest of the section, absent express mention to the contrary, we shall denote by X a locally compact space, by \(\Lambda: t \mapsto \lambda_{t}\) a \(\mu\) -adequate mapping of T into \(\mathcal{M}_{+}(X)\) , and by \(\nu\) the integral of \(\Lambda\) .

PROPOSITION 3. --- Let \(f\) be a numerical function \(\geqslant 0\) defined on X. a) The following inequalities hold:

\[
\int^ {*} f (x) d \nu (x) \geqslant \int^ {\bullet} d \mu (t) \int^ {*} f (x) d \lambda_ {t} (x) \geqslant \int^ {\bullet} d \mu (t) \int^ {\bullet} f (x) d \lambda_ {t} (x). \tag {6}
\]

\begin{enumerate}
\def\labelenumi{\alph{enumi})}
\setcounter{enumi}{1}
\tightlist
\item
  If \(\Lambda\) is vaguely continuous, then
\end{enumerate}

\[
\int^ {*} f (x) d \nu (x) \geqslant \int^ {*} d \mu (t) \int^ {*} f (x) d \lambda_ {t} (x).\tag{7}
\]

\begin{enumerate}
\def\labelenumi{\alph{enumi})}
\setcounter{enumi}{2}
\tightlist
\item
  If \(\lambda_t^\bullet (1) <   + \infty\) locally \(\mu\) -almost everywhere, then
\end{enumerate}

\[
\int^ {\bullet} f (x) d \nu (x) \geqslant \int^ {\bullet} d \mu (t) \int^ {*} f (x) d \lambda_ {t} (x) = \int^ {\bullet} d \mu (t) \int^ {\bullet} f (x) d \lambda_ {t} (x). \tag {8}
\]

Let \(g\) be a lower semi-continuous function on \(X\) such that \(f \leqslant g\) . For every \(t \in T\) ,

\[
\int^ {*} f (x) d \lambda_ {t} (x) \leqslant \int^ {*} g (x) d \lambda_ {t} (x),
\]

therefore, by (4) and Prop. 4 of §1,

\[
\int^ {\bullet} d \mu (t) \int^ {*} f (x) d \lambda_ {t} (x) \leqslant \int^ {\bullet} d \mu (t) \int^ {*} g (x) d \lambda_ {t} (x) = \int^ {*} g (x) d \nu (x).
\]

The first of the inequalities (6) then follows from the definition of \(\int^{*}f(x)d\nu(x)\) (Ch. IV, §1, No.~3, Def. 3), and the second follows immediately from it. The inequality (7) is proved in an analogous way if \(\Lambda\) is vaguely continuous, using (5) instead of (4).

Let us pass to the proof of (8). The mapping \(t \mapsto \lambda_t^*(1)\) is measurable, and finite locally \(\mu\) -almost everywhere. The set \(\mathfrak{K}\) of compact subsets of \(T\) such the restriction of \(t \mapsto \lambda_t^*(1)\) to \(K\) is finite and continuous is therefore \(\mu\) -dense, and Prop. 4 of §2, No.~3 implies the existence of a summable family \((\mu_\alpha)_{\alpha \in A}\) of positive measures, with supports belonging to \(\mathfrak{K}\) , such that \(\mu = \sum_{\alpha \in A} \mu_\alpha\) . The mapping \(\Lambda\) is \(\mu_\alpha\) -adequate for every \(\alpha \in A\) ; set \(\nu_\alpha = \int \lambda_t d\mu_\alpha(t)\) . Prop. 1 shows that \(\nu = \sum_{\alpha \in A} \nu_\alpha\) , and the relation (4), applied to the measure \(\mu_\alpha\) and the function 1, shows that \(\nu_\alpha\) is a bounded measure (because \(\lambda_t^*(1)\) is bounded on \(\operatorname{Supp}(\mu_\alpha)\) ). Let us then write the formula (6) for the measure \(\mu_\alpha\) , replacing the symbol \(\int^*\) in the first member by \(\int^\bullet\) , which is legitimate by Prop. 7 of §1; then

\[
\int^ {\bullet} f (x) d \nu_ {\alpha} (x) \geqslant \int^ {\bullet} d \mu_ {\alpha} (t) \int^ {*} f (x) d \lambda_ {t} (x) = \int^ {\bullet} d \mu_ {\alpha} (t) \int^ {\bullet} f (x) d \lambda_ {t} (x)
\]

(the last equality due to the fact that \(\lambda_t\) is bounded locally almost everywhere, and Prop. 7 of § 1). The inequality in (8) is then obtained by summing on \(\alpha\) (§2, No.~2, Prop. 1).

If no hypothesis analogous to that of c) is made, the inequality (8) may fail (Exer. 2).

COROLLARY 1. --- Let \(f\) be a function \(\geqslant 0\) defined on \(X\) , and let \(H\) be the set of \(t \in T\) such that \(f\) is not \(\lambda_t\) -negligible.

\begin{enumerate}
\def\labelenumi{\alph{enumi})}
\item
  If \(f \nu\) -negligible, then \(H\) is locally \(\mu\) -negligible.
\item
  If \(f\) is \(\nu\) -negligible and \(\Lambda\) is vaguely continuous, then \(H\) is \(\mu\) -negligible.
\item
  If \(f\) is locally \(\nu\) -negligible and \(\lambda_t^\bullet(1) < +\infty\) locally \(\mu\) -almost everywhere, then \(H\) is locally \(\mu\) -negligible.
\end{enumerate}

COROLLARY 2. --- Let f be a function \(\geqslant0\) defined on X, \(\nu\) -measurable and \(\nu\) -moderated. The set of \(t\in T\) such that f is not \(\lambda_{t}\) -moderated is then locally \(\mu\) -negligible (and even \(\mu\) -negligible if \(\Lambda\) is vaguely continuous).

For, \(f\) is the sum of a sequence of functions \(f_{n} \geqslant 0\) such that \(f_{n}\) is zero outside a compact set \(K_{n}\) for \(n \geqslant 1\) , and \(f_{0}\) is \(\nu\) -negligible (§1, No.~2, Prop. 6); \(f_{0}\) is then \(\lambda_{t}\) -negligible except for \(t\) forming a set that is locally \(\mu\) -negligible (and even \(\mu\) -negligible, if \(\Lambda\) is vaguely continuous) by Cor. 1, and the statement then follows at once.

PROPOSITION 4. --- Let f be a \(\nu\) -measurable function defined on X, with values in a topological space G, and let M be the set of \(t \in T\) such that f is not \(\lambda_{t}\) -measurable.

\begin{enumerate}
\def\labelenumi{\alph{enumi})}
\item
  Suppose that \(f\) is constant on the complement of a \(\nu\) -moderated subset of \(X\) ; then \(M\) is locally \(\mu\) -negligible.
\item
  Suppose that \(f\) is constant on the complement of a \(\nu\) -moderated subset of \(X\) , and that \(\Lambda\) is vaguely continuous; then \(M\) is \(\mu\) -negligible.
\item
  Suppose that \(\lambda_t^\bullet(1) < +\infty\) locally \(\mu\) -almost everywhere; then M is locally \(\mu\) -negligible.
\end{enumerate}

Let us first prove a) (resp. b)). Since every \(\nu\) -integrable set is contained in a \(\nu\) -integrable open set, the function f is constant on the complement B of a countable union of \(\nu\) -integrable open sets. There exists a partition of X - B formed by a \(\nu\) -negligible set N and a sequence \((\mathrm{K}_{n})\) of compact sets such that the restriction of f to each \(K_{n}\) is continuous. Let S be the set of \(t \in T\) such that N is not \(\lambda_{t}\) -negligible: S is locally \(\mu\) -negligible (resp. \(\mu\) -negligible) by Cor. 1 of Prop. 3. The sets \(K_{n}, B, N\) are measurable for every measure on X, and the restriction of f to each of them is \(\lambda_{t}\) -measurable for every \(t \notin S\) . The function f is therefore \(\lambda_{t}\) -measurable for every \(t \notin S\) (Ch. IV, §5, No.~10, Prop. 16).

To establish c), let us take up again the notations in the proof of Prop. 3; since f is \(\nu\) -measurable, it is measurable for each of the measures \(\nu_{\alpha} \leqslant \nu\) . Now, these measures are bounded, hence moderated, and it follows from a) that M is locally \(\mu_{\alpha}\) -negligible for every \(\alpha \in A\) . This implies that M is locally \(\mu\) -negligible (§2, No.~2, Cor. 2 of Prop. 1).

PROPOSITION 5. --- Let f be a \(\nu\) -measurable positive numerical function defined on X, and let N be the set of \(t \in T\) such that f is not both \(\lambda_{t}\) -measurable and \(\lambda_{t}\) -moderated.

\begin{enumerate}
\def\labelenumi{\alph{enumi})}
\tightlist
\item
  Suppose that \(f\) is \(\nu\) -moderated. The set \(\mathbf{N}\) is then locally \(\mu\) -negligible, the function \(t \mapsto \int^{\bullet} f(x) d\lambda_t(x)\) is \(\mu\) -measurable, and
\end{enumerate}

\[
\int^ {\bullet} f (x) d \nu (x) = \int^ {\bullet} d \mu (t) \int^ {\bullet} f (x) d \lambda_ {t} (x).\tag{9}
\]

\begin{enumerate}
\def\labelenumi{\alph{enumi})}
\setcounter{enumi}{1}
\tightlist
\item
  Suppose that \(f\) is \(\nu\) -moderated, and that \(\Lambda\) is vaguely continuous. The set \(N\) is then \(\mu\) -negligible, the function \(t \mapsto \int^{*} f(x) d\lambda_{t}(x)\) is \(\mu\) -measurable and \(\mu\) -moderated, and
\end{enumerate}

\[
\int^ {*} f (x) d \nu (x) = \int^ {*} d \mu (t) \int^ {*} f (x) d \lambda_ {t} (x).\tag{10}
\]

\begin{enumerate}
\def\labelenumi{\alph{enumi})}
\setcounter{enumi}{2}
\tightlist
\item
  Suppose that \(\lambda_t^\bullet(1) < +\infty\) locally \(\mu\) -almost everywhere. The set N is then locally \(\mu\) -negligible, the function \(t \mapsto \int^\bullet f(x) d\lambda_t(x)\) is \(\mu\) -measurable, and (9) holds.
\end{enumerate}

Let us first prove a) (resp. b)), assuming that \(f\) is \(\nu\) -moderated. The assertions concerning the set \(N\) have already been established (Prop. 4, and Cor. 2 of Prop. 3). By Prop. 6 of §1, No.~2, we may limit ourselves to proving a) (resp. b)) in each of the following special cases:

\begin{enumerate}
\def\labelenumi{\arabic{enumi})}
\item
  The function \(f\) is \(\nu\) -negligible.
\item
  There exists a compact set \(\mathbf{K}\) such that \(f\) is zero outside \(\mathbf{K}\) and the restriction of \(f\) to \(\mathbf{K}\) is continuous.
\end{enumerate}

The special case 1) has already been treated (Cor. 1 of Prop. 3). To treat the second, we denote by G a \(\nu\) -integrable open set containing K, by M a constant upper bound for f, by h the lower semi-continuous function \(M\varphi_{G}\) , and by g the function h-f.~Since the function f is upper semi-continuous on X, g is lower semi-continuous and positive. Moreover, f,g,h are \(\nu\) -integrable.

Let us then apply formula (4) (resp. (5)) to the lower semi-continuous functions h and g. By subtraction, we see that the function

\[
t \mapsto \int^ {\bullet} f (x) d \lambda_ {t} (x) \quad (\text { resp. } \int^ {*} f (x) d \lambda_ {t} (x))
\]

is \(\mu\) -measurable and that the formula (9) (resp. (10)) holds. Finally, under the hypothesis of b), the function \(t \mapsto \int^{*} f(x) d\lambda_{t}(x)\) has finite upper integral, hence is indeed \(\mu\) -moderated.

To prove c), let us take up again the measures \(\mu_{\alpha}\) and \(\nu_{\alpha}\) of the proof of Prop. 3; since f is \(\nu_{\alpha}\) -measurable and \(\nu_{\alpha}\) -moderated, the assertion a) implies that \(t \mapsto \int^{\bullet} f(x) d\lambda_{t}(x)\) is \(\mu_{\alpha}\) -measurable and that

\[
\int^ {\bullet} f (x) d \nu_ {\alpha} (x) = \int^ {\bullet} d \mu_ {\alpha} (t) \int^ {\bullet} f (x) d \lambda_ {t} (x).
\]

It remains only to sum on \(\alpha\) , applying Props. 1 and 2 of §2, No.~2.

If \(f\) is not assumed to be \(\nu\) -moderated, and if one does not make the assumption in c), then the relation (9) may not hold (Exer. 3).

COROLLARY. --- Let f be a function defined on X, with values in a Banach space F or in \(\overline{R}\) , that is \(\nu\) -measurable and \(\nu\) -moderated. For f to be \(\nu\) -integrable, it is necessary and sufficient that

\[
\int^ {\bullet} d \mu (t) \int^ {\bullet} | \mathbf {f} (x) | d \lambda_ {t} (x) <   + \infty .
\]

This follows at once from Prop. 5 and the criterion for integrability (Ch. IV, §5, No.~6, Th. 5).

